\section{Literature}\label{sec:literature}

This paper sits at the intersection of five literatures: the incidence of trade shocks; the evidence on how workers adjust to them; the institutional analyses of the 2025 US tariffs on the UK; the public-finance literature on progressive taxation as social insurance, which supplies the mechanism; and the microsimulation literature on tax--benefit cushioning of earnings shocks, which supplies the method. Existing macro--micro studies impose labour-market shocks on tax--benefit models, while tariff studies map price and income exposure into heterogeneous household welfare---the World Bank's Household Impacts of Tariffs framework combines household consumption and income portfolios \citep{artuc2021}, and the OECD links its METRO trade model to household-budget surveys \citep{luu2020}. The narrower contribution here is a UK application that maps the 2025 partner-country export-demand exposure into alternative labour-income stress scenarios and then through the statutory tax--benefit system; relative to those economy-wide tools it supplies only the terminal tax--benefit module, conditioning on an imposed labour-income distribution rather than deriving it from production (see Section~\ref{sec:method-primitives}).

\subsection{Trade shocks and their incidence}

Two recent papers anchor the structural end and frame the adjustment-margin question as its poles. \citet{ignatenko2025} build a quantitative general-equilibrium model of the 2025 tariffs in which wages are flexible and labour supply elastic, so employment is an equilibrium outcome whose sign flips with revenue recycling. \citet{atkin2025} instead apply a wedge framework to Chilean administrative data with factor supplies fixed: all adjustment runs through flexible wages and prices. Their object is the incidence of domestic distortions---markups, taxes, credit frictions---rather than of tariffs, and their vocabulary transfers here only in part. They distinguish the ``statutory'' incidence of a wedge, meaning who hands over the money given fixed prices, from the ``economic'' incidence that emerges once factor prices adjust in general equilibrium, and their contribution is to compute the second. This paper computes only the first: factor prices are held fixed by assumption throughout (Section~\ref{sec:method-primitives}), so what is measured is the statutory incidence of an imposed earnings loss, and the economic incidence is ceded to the general-equilibrium papers on either side of it. I borrow their distinction as a way of naming what this design does and does not deliver, not as the framework I implement. Between the poles, \citet{rodriguezclare2025} embed the 2025 trade war in a model with downward nominal wage rigidity, in which rigid wages turn the same tariff primitive into an employment shock that flexible-wage models turn into a wage shock---precisely the disagreement my adjustment-margin axis makes an explicit scenario dimension. \citet{michelena2025} find 2025 employment and export losses concentrated in the targeted countries, and \citet{denbestenkanzig2025} show, in a narrative series over 1840--2024, that tariff increases are robustly contractionary---discipline for treating the shock as a negative export-demand shock.

The reduced-form lineage begins with \citet{autor2013}: import competition caused large, persistent employment and wage losses in exposed US commuting zones \citep[effects persisting to 2019;][]{autor2021}. \citet{caliendo2019} provide the structural counterpart, \citet{adao2023} the mapping from regional estimates to aggregates, \citet{kimvogel2021} the money-metric worker welfare translation with regressive incidence, and \citet{galle2023} the coexistence of aggregate gains with group-specific losses. My contribution relative to this work is to replace stylised worker groups and transfers with real households and the actual UK tax--benefit system. The 2018--19 US episode supplies the sharpest evidence on who pays a tariff: pass-through into US import prices was essentially complete \citep{amiti2019, fajgelbaum2020}, retail prices adjusted more slowly \citep{cavallo2021, cavallo2025}, and US manufacturing employment fell on net \citep{flaaenpierce2019}. The retaliation side is the closest analogue to my channel---a partner's tariff arriving as an export-demand shock \citep{waugh2019, blanchard2019}. The broader distributional literature establishes redistribution along both consumption and income margins \citep{fajgelbaumkhandelwal2016, fajgelbaumkhandelwal2022, borusyakjaravel2021}. Two propagation channels my direct-channel design excludes have well-measured analogues: supply-chain transmission \citep{barrotsauvagnat2016, acemoglu2016} and local-services spillovers \citep{moretti2010, gathmann2020}; Section~\ref{sec:discussion} carries the caveat. Exchange rates absorb part of a tariff, imperfectly and state-dependently \citep{jeanneson2024, corsettilloydostry2025, kalemliozcan2026, auray2020}; I hold the exchange rate fixed. Early assessments of the 2025 wave again find the incidence landing on the tariff-imposing country's consumers \citep{fajgelbaumkhandelwal2026}---complementary to my question, the incidence on the tariffed \emph{partner's} workers.

\subsection{Adjustment margins: how workers absorb trade shocks}

The worker-level evidence says the margin is mostly earnings, not permanent unemployment. \citet{autor2014} find exposed US workers lose earnings amid churn rather than indefinite joblessness; \citet{traiberman2019} shows occupation-specific human capital produces reallocation with persistent wage losses. Because my shock is an export-demand shock rather than import competition, the most direct precedent is \citet{garinsilverio2024}, who use Portuguese matched employer--employee data to split a firm's export-demand shock into a product-market component and an idiosyncratic component, and find that 10--15 per cent of the idiosyncratic component passes through into incumbent employees' wage growth, with no effect on retention rates. The magnitude should be quoted as it is estimated: it is a pass-through range for an \emph{idiosyncratic firm-level} demand shock in a recession, not a point elasticity of wages to a sector-level exposure index, and the absence of a retention response makes it evidence on the intensive margin specifically. It is the sharpest available warrant for including a wage-cut dimension at all, and---because the pure wage-cut family delivers the whole sector wage-bill loss through incumbent pay, that is, at unit pass-through---it is also the clearest evidence that that family is an upper bound rather than a central case (Section~\ref{sec:method-margins}). \citet{hummels2014} is a companion rather than a precedent: it estimates the wage effects of \emph{offshoring} in Danish matched worker--firm data, finding economically substantial within-worker effects differentiated by skill, on a design that does not identify export-demand pass-through. The rent-sharing literature disciplines its plausible scale: a survey of twenty-two studies placing elasticities of wages to firm-specific rents at 0.05--0.15 \citep{cardcardosoheining2018}; \citet{klinepetkova2019} put the worker share of patent-induced surplus at roughly 30 cents in the dollar, a related but distinct object. \citet{dauth2021} provide the cleanest evidence for the reallocation family, \citet{dixcarneiro2014} the adjustment-cost horizons and age gradients, with \citet{dixcarneirokovak2017} showing that the regional consequences of those frictions keep widening for two decades, and \citet{rud2022} decompose displacement costs into wage cuts on re-employment versus non-employment spells---directly calibrating the split between my families.

The UK's own history points to a third margin: benefit-financed economic inactivity. \citet{beatty2007} document that coalfield job losses became male inactivity on incapacity benefits, and \citet{beatty2017} show that 1980s--90s industrial job destruction still inflates the deficit---the direct UK precedent for my fiscal-incidence question, retrospective and area-level where mine is prospective and household-level. The margin is live: disability claims still respond to local slack \citep{robertstaylor2022}, hidden unemployment persists \citep{beatty2023}, and working-age health-related benefit rolls in Great Britain reached 4.2 million by August 2023 \citep{ifs2024}. For UK-specific elasticities, \citet{delyonpessoa2021} show workers in Chinese-import-competing industries lost employment years and earnings \citep[see also][]{dorn2024, aragon2018}. Closest to my object, \citet{irastorza2025} study UK household responses to trade shocks---added-worker effects, retirement timing, self-employment---but model behaviour, not tax--benefit mechanics; I treat their margins as complementary.

\subsection{The 2025 tariffs and the UK}

The institutional literature on the 2025 shock is rich but stops above the household. The OBR's March 2025 scenarios put the permanent UK GDP loss from a 20-point US tariff at around 0.3 per cent, rising to about three-quarters of a per cent in the medium term only under universal retaliation; a third scenario with the same US rise leaves UK GDP largely unchanged \citep{obr2025}; NIESR provides NiGEM scenarios \citep{niesr2024}. The Bank of England estimates the effective US tariff on the UK at roughly 9 per cent and finds trade diversion has mildly \emph{lowered} UK import prices \citep{boe2025, dhingra2025}---which licenses isolating the export-demand/earnings channel while holding price and monetary channels fixed. The IMF and WTO frame the aggregates \citep{imf2025, wto2025}; UK firm surveys find mostly small expected effects, with the US taking about 22 per cent of UK exports \citep{cepdmp2025}. The geography is sharp: Coventry sends 22 per cent of its exports to the US \citep{centreforcities2025}; auto tariffs are estimated to cost \pounds 10,000 million of GDP over five years, 62 per cent borne by the West Midlands \citep{cityredi2025}; SMMT data show US vehicle exports collapsing between April and May 2025 \citep{smmt2025}. Some 80 per cent of UK steel exports go to Europe---the basis of my steel caveat---and the December 2025 pharmaceutical deal was paid for through VPAG rebate cuts and NHS pricing concessions, making pharma a fiscal/pricing shock rather than an employment shock \citep{abpi2025}. The Resolution Foundation's household-facing analysis is descriptive rather than microsimulated \citep{resolutionfoundation2025}, and the IFS has no dedicated distributional publication on the 2025 tariffs; that gap defines the contribution. For UK trade-shock precedents, \citet{dhingra2017} map Brexit trade costs to incomes, \citet{fetzerwang2024} provide district-level incidence, \citet{fetzer2019} and \citet{bloom2019} the political-economy and uncertainty channels, and \citet{breinlich2022} and \citet{levell2017} benchmark the consumption-price channel that this cash-income exercise holds fixed.

\subsection{Progressive taxation as social insurance: the marginal--average wedge}\label{sec:lit-progressivity}

Asking a tax--benefit system to absorb a trade shock is not an arbitrary pairing of literatures. \citet{rodrik1998} documents the compensation hypothesis---more open economies systematically run larger governments, a correlation he attributes to the public sector insuring citizens against the external risk that openness brings---which is the reason a tariff and a tax schedule belong in the same paper at all. \citet{antras2017} make the welfare consequence precise: when redistribution is costly, the gains from globalisation depend on the instruments available to redistribute them, and an economy whose instruments are weak can lose from trade that it would gain from otherwise. Closest of all is \citet{lyonwaugh2018}, who put a progressive income tax inside an incomplete-markets trade model in which import competition raises uninsurable worker-level risk, and find that optimal progressivity should rise with openness---roughly five percentage points on top marginal rates for a ten-point increase in openness. Their exercise and mine are the two halves of one question. They solve for the progressivity a trading economy should choose; I take the progressivity the UK has actually legislated as given and ask how much of a specific aggregate earnings loss it absorbs, and how that answer moves when the same aggregate is redistributed across workers. Neither question is answerable without the other's object.

That object is the wedge between marginal and average tax rates. The idea that a progressive schedule is an insurance contract rather than only a redistributive one originates with \citet{varian1980} and \citet{eatonrosen1980}: when income contains an uninsurable random component, a positive marginal rate is optimal as risk-sharing quite apart from any social taste for equality, because the state takes a share of the bad draw as well as of the good one. \citet{heathcote2017} turn that idea into a measurable parameter, and their framework is the one this paper's mechanism is a statement about. Under the log-linear schedule $T(y)=y-\lambda y^{1-\tau}$, one minus the marginal rate equals $(1-\tau)$ times one minus the average rate, so $m-a=\tau(1-a)$: the wedge is $\tau$ scaled by one minus the average rate, and equals $\tau$ only where the average rate is zero, and they interpret it directly as the insurance content of the schedule---the share of an idiosyncratic earnings shock the tax system absorbs. My central mechanism is that wedge restated in levels: a partial earnings cut is relieved at the marginal rate, a complete earnings loss at the average rate, so the same aggregate loss is cushioned less when it is concentrated on fewer workers. The cushioning rate reported here is therefore an estimate of the effective marginal--average wedge of the UK tax--benefit schedule, evaluated at the exposed-earnings distribution rather than fitted as a single $\tau$ for the population. Two consequences shape the design. First, the UK schedule is emphatically not log-linear---the personal allowance, the National Insurance thresholds, the Universal Credit taper and the discrete arrival of out-of-work entitlement at zero earnings all make the local wedge jump with the level and the completeness of the loss---so no single $\tau$ summarises it and the exercise has to be simulated on the actual rules rather than solved. Second, in the \citet{heathcote2017} framework the wedge is a property of the schedule alone, whereas here it is a property of the schedule \emph{and} of who is shocked; that is precisely why redistributing a fixed aggregate across the exposed population moves the measured number, and it is the sense in which this paper's result is an addition to theirs rather than an application of it.

The empirical counterpart to the wedge is the partial-insurance literature, which recovers the same quantity from data instead of from a schedule. \citet{blundellpistaferripreston2008} estimate what share of a transitory or permanent earnings shock passes into consumption, finding transitory shocks close to fully insured and permanent shocks only partially so; \citet{kaplanviolante2010} ask how much of that measured insurance a life-cycle model with self-insurance through saving alone can reproduce, and how much requires something more; and \citet{blundellgrabermogstad2015}, in Norwegian administrative data, decompose what remains into the part supplied by taxes and transfers and the part supplied by the family. Their estimand is this paper's estimand with an empirical rather than a simulated answer: the share of an earnings shock that does not reach household resources. Two differences define the division of labour. They measure pass-through into \emph{consumption}, so their insurance includes saving, borrowing and family labour supply, whereas I stop at disposable income and therefore isolate the tax--benefit component that \citet{blundellgrabermogstad2015} identify as one term among several. And they recover an average over the shocks a population happens to experience, whereas the counterfactual here holds the aggregate fixed and varies its distribution---the one dimension a pass-through regression estimated on realised shocks cannot hold fixed. Their finding that insurance is markedly weaker against large and persistent shocks than against small transitory ones is the closest independent corroboration of the direction of this paper's result, and a standing caution against reading any statutory replacement rate as the whole of what a household is protected by.

\subsection{Tax--benefit cushioning of earnings shocks}

The methodological ancestors are \citet{dolls2012}, who show EU tax--benefit systems absorb about 38 per cent of a proportional income shock and 47 per cent of an unemployment shock. Their \emph{UK} cells are the right comparator: a UK income-shock stabilisation coefficient of 0.352 (income tax 0.267, social insurance 0.054, benefits 0.031) and a UK unemployment-shock coefficient of 0.415 (tax 0.191, SIC 0.061, out-of-work benefits 0.163). My \SubmissionUnitWageCushion\ per cent wage-cut and $\SubmissionUnitDisplacedCushion \pm \SubmissionUnitDisplacedCushionSD$ per cent displacement rates are of a similar order, but the \emph{ordering is reversed}: Dolls et al.\ find the unemployment shock more strongly stabilised than the income shock, and I find the opposite. That reversal is the sharpest external check available on this paper's headline contrast, and Section~\ref{sec:results-benchmarks} reconciles the two decompositions channel by channel rather than leaving the discrepancy to the reader. The marginal-rate framing of the income tax as an automatic stabiliser originates with \citet{auerbachfeenberg2000} and is formalised by \citet{mckayreis2016}; \citet{kniesnerziliak2002} supply the direct evidence that the wedge of Section~\ref{sec:lit-progressivity} is a policy choice rather than a constant, showing that the US tax reforms of the 1980s, by flattening the schedule, measurably reduced the consumption insurance the income tax provided. Their result is the historical warrant for treating a measured cushioning rate as something a legislature moves, which is what makes the counterfactual schedules in Section~\ref{sec:policy} more than an arithmetic exercise. \citet{figari2011} establish EUROMOD's stress-test use, on the model documented by \citet{sutherlandfigari2013}; \citet{fernandezsalgado2014} give the closest published estimate of how far European welfare systems compensated households for job loss in the Great Recession, and are the natural external comparator for the displacement margin here; my direct methodological sibling is \citet{brewertasseva2021}, who pass the Covid-19 labour-market shock through UKMOD---the same architecture, applied here to a trade-policy shock without the emergency schemes and with the adjustment margin as the central axis. The macro--micro linkage literature supplies the shock-derivation lineage \citep{bourguignon2008, peichl2016, navicke2013}. On the US side, exposed commuting zones saw rising disability, retirement and welfare payments offsetting only a small fraction of losses \citep{autor2013}, and \citet{hyman2024} show wage insurance raises re-employment---a candidate counterfactual reform for this framework. The nearest tariff-microsimulation links are US consumption-side models \citep{tpc2025, budgetlab2025, acostacox2024}. Two institutional facts motivate the cushioning asymmetry examined here: UC's capital rules---thresholds frozen in cash terms since 2006---reduced or extinguished entitlement for around two million income-eligible families \citep{resolutionfoundation2025saving}, and take-up of income-related benefits is well below 100 per cent \citep{dwp2024takeup}, so statutory replacement rates can overstate what displaced workers receive. Take-up is examined as a headline sensitivity (Sections~\ref{sec:method-takeup} and~\ref{sec:results-takeup}); the capital rules are a real-world gate the plain FRS data cannot exercise (Section~\ref{sec:discussion}).

Because claiming rather than entitlement is the binding parameter in the displacement scenarios, the take-up evidence deserves more weight than the official series alone can carry. \citet{currie2006} surveys take-up across programmes and finds it systematically incomplete, and systematically lowest where entitlements are small and where the transaction and stigma costs of claiming are large. That gradient matters here in a specific direction: a household that has just lost its main earnings has an unusually large entitlement and so sits at the high-take-up end of it, which is why the take-up sensitivity should be read as a bound on how far incomplete claiming can erode measured cushioning rather than as a central estimate. \citet{bhargavamanoli2015} sharpen the point with an IRS field experiment in which a substantial share of non-claiming among people already notified of their eligibility for the EITC responds to the complexity and presentation of the claim itself rather than to information about eligibility. Take-up is therefore an artefact of administrative design rather than a fixed parameter of the population, which is what licenses treating the take-up scenarios of Section~\ref{sec:method-takeup} as a policy dimension of the cushioning result and not merely as a robustness check on it.

What a pound of benefit is worth to the household receiving it is a further step this cash-income design does not take, and two papers mark the size of the gap. \citet{ganongnoel2019}, using de-identified bank-account data, find non-durable spending falls by about 6 per cent at the onset of unemployment and by a further 12 per cent at the exhaustion of unemployment insurance, and conclude that the consumption-smoothing gains from extending benefit duration are roughly four times those from raising benefit levels---a finding that bears directly on the duration dimension of the scenarios here, in which the same annual earnings loss is delivered over three months or over twelve. \citet{landaisspinnewijn2021} estimate the value of that insurance directly in Swedish data and put the marginal value of a transfer at least 60 per cent higher when unemployed than when employed, well above what the relatively modest consumption drop would imply on its own. The cushioning rate is silent on both results: it counts pounds of disposable income, holds saving, borrowing and the timing of expenditure fixed, and treats a pound of benefit and a pound of earnings as interchangeable. The direction of the resulting understatement is worth stating plainly---a cash measure records none of the extra welfare value that benefit income carries precisely in the state where it arrives, which is the displacement margin---and Section~\ref{sec:discussion} carries the caveat.
